\section{Census input and commutativity}\label{sec:data}
The input is the pinned \texttt{fusion-ring-census} release
\cite{Census}. It counts based-ring isomorphism classes, rather than
raw multiplication tables; its exact-multiplicity counts through twenty
are displayed in Table~\ref{tab:census}. Our exclusion ledger addresses
each rank-five tensor through eighteen. To connect this archived ledger
with the repository tables, we canonically flatten every tensor under
all $24$ permutations of the four nonunit basis elements. The resulting
sets agree at each multiplicity $1\le\mu\le18$: all $28493$ repository
tensors have an isomorphic archived tensor, and conversely. The
multiplicity-nineteen and -twenty certificates address the repository
tensors directly, in their recorded repository order. Labels do
not generally agree across these releases, so the classification
transport uses based-ring isomorphism, never equality of row indices
or of labeled-tensor fingerprints. The census identifies the input;
the categorical conditions and certificates decide its members.

The following argument supplies the commutativity needed for both the
ordinary and signed sphericalization computations. The unsigned assertion
also appears in \cite[Proposition 9.2]{ABDP}.

\begin{lemma}\label{lem:commutativity}
Every rank-five fusion ring is commutative. A rank-five signed Frobenius
based ring admitting a nonzero, dual-equal dimension character is also
commutative.
\end{lemma}
\begin{proof}
The complex algebra of a (possibly signed) Frobenius based ring is
semisimple: the coefficient-of-unit functional gives the positive
inner product making the basis orthonormal and multiplication a
$*$-representation. A noncommutative algebra of dimension five with a
one-dimensional representation must be $\CC\oplus M_2(\CC)$.
Its unique one-dimensional character is Galois invariant, hence rational,
and its values are algebraic integers, hence integers.
For a fusion ring this character is FP; in the signed assertion it is
the given dimension character. In either case its formal codegree is an
integer $D\geq5$, the sum of five nonzero integer squares.
If $f$ denotes the formal codegree of the two-dimensional representation,
the standard identity $\sum_E\dim(E)/f_E=1$ gives
\[
 f=\frac{2D}{D-1}=2+\frac2{D-1}.
\]
Formal codegrees are algebraic integers. The displayed rational number
is not an integer for $D\geq5$, a contradiction.
\end{proof}

In the self-dual case, our tensor coordinate convention is the twenty
ordered nonunit triples
\[
111,112,113,114,122,123,124,133,134,144,
222,223,224,233,234,244,333,334,344,444.
\]
The coefficients $v_j$ are listed in this order; multiplication with the
unit is fixed by the based-ring axioms. The two tensors resolved by
inert-prime parity in
Section~\ref{sec:automatic} are recorded in these coordinates. For
nonself-dual inputs, the data contain the full multiplication tensors.

\begin{table}[ht]\centering\small
\begin{tabular}{rrrr}\toprule
$\mu$&Number of rings&$\mu$&Number of rings\\\midrule
1&16&11&1383\\
2&37&12&1948\\
3&82&13&2211\\
4&134&14&2554\\
5&209&15&2936\\
6&336&16&4157\\
7&477&17&4255\\
8&733&18&5080\\
9&863&19&5640\\
10&1082&20&6866\\
\bottomrule\end{tabular}
\caption{Rank-five based-isomorphism classes at exact multiplicity, from the pinned census \cite{Census}. The categorification decisions in this paper cover every displayed multiplicity.}\label{tab:census}
\end{table}

